\section{Preliminaries for differentially private data synthesis}
In this section, we summarize key concepts and introduce notation related to differentially private data synthesis used in the paper. 

In general, we consider our data to lie in a domain $\domain= \domain_1 \times  \cdots \times \domain_d$ that is discrete and has dimension $d$.
This assumption is not restrictive since the vast majority of DP-data synthesis algorithms for tabular data rely on a discretized version of the data even if it originally lies in a continuous domain\footnote{We refer the reader to \citep{zhang2016privtree} for a discussion on how to optimally discretize in a DP-way.}. In particular, we can represent every dimension as an integer in the discrete set $\domain_{i} = \{1, \cdots, k_i\}$ with $k_i \in \mathbb N_+$.

%\subsection{Differentially private data synthesis}
%\fy{section title doesn't capture section - neither discretization nor DP definition are about dp data synthesis}
Differential privacy \citep{dp} is an algorithmic property that guarantees that individual information in the data is protected in the output of an algorithm; even when assuming that an adversary has access to the information of all other individuals in the dataset. We now provide the formal definition.
%\fy{ this belongs to definition} More formally, for any dataset $\dataset \in \domain^n$ of size $n$ in some domain $\domain$, 
%we define: \fy{grammatically this has to be correct together with definition, but this is not}

\begin{definition}
\label{eq:defdp}
    An algorithm $\gA$ is $(\epsilon,\delta)$-DP with $\epsilon>0$ and $\delta>0$ if for any datasets $\dataset, D'$
    %$\dataset \in \domain^n$ of size $n$
    differing in a single entry and any measurable subset $S \subset \mathrm{im}(\mathcal A)$ of the image of $\mathcal A$, we have
\begin{equation*}
  \prob{\algo{D}\in S} \leq \exp(\epsilon) \prob{\algo{D'}\in S} + \delta
\end{equation*}
\end{definition}

The goal of \emph{differentially private data synthesis} is to design a (randomized) algorithm $\gA$  that, for any dataset $\dataset \in \domain^n$ of size $n$, generates an output $\algo{D} \in \domain^m$ that is a differentially private ``copy" of $D$, potentially of different size $m \neq n$. 

% \paragraph{Discretization}  

%While the original dataset may contain both continuous and categorical entries, we consider a discretized version of the data in a domain $\domain = \domain_1 \times  \cdots \times \domain_d$ with dimension $d$ with $\domain_{i} = \{1, \cdots, k_i\}$ in \fy{in?} at most $ k_i \leq k$ bins per dimension. 
%Discretization is a crucial ingredient required by the vast majority of DP-data synthesis algorithms designed for tabular data, such as marginal-based approaches (see next subsection).  
%There exists \fy{"there exists" generally a phrase to avoid except mathematically} a separate line of works dedicated to finding optimal discretizations in a DP-way, which is orthogonal to the work presented in this paper; see, for example, \citep{zhang2016privtree}. 

\subsection{Marginal-based algorithms for private data synthesis }
\label{sec:marginal_based}
Our approach falls in the general category of marginal-based methods. They
%Marginal-based approaches 
 follow Algorithm~\ref{alg:mainalgo},  consisting of three steps: marginal selection, privatization, and generation.

\paragraph{Marginal selection} 
For any subset $S \subset \{1, \cdots, d\}$ of the dimensions,  we denote with $D_S \in \domain_S^n$ the dataset containing only the dimensions in $S$. For each subset $S$ we can define a corresponding marginal.
\begin{definition}
\label{def:marginal}
    We denote by $\marginal{S}[D] \in \mathcal P(\domain_S)$ the $S$-marginal of a dataset $D$, defined as the empirical measure of $D_S$ over the domain $\domain_S$.
    \end{definition}

In the first step of the algorithm, a set $\mathcal{S}$ of such subsets $S$ is selected, and equivalently a set of marginals. 
%subsets $S \in \mathcal S$ \fy{is S a subset or set of subsets}: \fy{youre not defining a set of subsets...}
%The set of subsets $\mathcal S$ defines the marginals. 
The problem of selecting marginals in a DP-way is an interesting problem on its own and has led to a significant amount of proposed methods~\citep{cai2021data, mst, aim, zhang2021privsyn}. 
%Thus, marginal selection methods are orthogonal to this paper.
In an
%One important 
extension of Algorithm~\ref{alg:mainalgo}, sketched in Algorithm~\ref{alg:mwem}, the marginals are not all selected in the beginning, but sequentially chosen from 
%is to sequentially select the marginals from 
a pre-defined pool of subsets $\mathcal S_W$ of $[d]$, often referred to as the workload.
%using the DP-dataset from the previous iteration. 
More specifically, in every iteration $t$, we select the marginals with the largest total variation distance $\text{TV}\left( \marginal{S}[\dpdata^{(t-1)}], \marginal{S}[D]\right) $, where $\dpdata^{(t-1)}$ is the DP-dataset from the previous iteration $t-1$. This is done
in a DP-way using the exponential mechanism \cite{mcsherry2007mechanism}. Such approaches fall under the general MWEM~\citep{mwem,liu2021iterative} framework, which is the backbone of many DP synthesis methods. To control the overall privacy budget, the framework from Algorithm~\ref{alg:mwem} uses advanced compositional theorems. 


\begin{figure}[!t]
\begin{minipage}[t]{.46\textwidth}
\noindent
\begin{algorithm}[H]
\caption{Standard data synthesis framework}\label{alg:mainalgo}
\begin{algorithmic}[1]
\Require Dataset $D$, privacy parameters $\epsilon$ and $\delta$
\State \textbf{select} set $\mathcal S$ of subsets $S$ of $\{1, \cdots, d\}$
\State \textbf{privatize} marginals $\marginal{S}[D]$ to obtain $(\epsilon, \delta)$-DP ``copies" $\measuredmarginal{S}$
\State \textbf{generate} data from privatized marginals $\measuredmarginal{S}$
\Statex\Return the DP dataset $\dpdata$
\end{algorithmic}
\end{algorithm}
\end{minipage}
\hfill
\begin{minipage}[t]{.52\textwidth}
\begin{algorithm}[H]
\caption{Extension to sequential query selection}\label{alg:mwem}
\begin{algorithmic}[1]
\Require Dataset \(D\), privacy parameters \(\epsilon\) and \(\delta\), workload \(\mathcal{S}_W\), rounds \(T\)
\For{\(t = 1, \dots, T\)}
    \State \textbf{select} \(S_t \in \mathcal{S}_W\)
    \State \textbf{privatize} the marginal \(\nu_{S_t}[D]\) to obtain the 
    \Statex\hspace{\algorithmicindent}\hspace{\algorithmicindent}DP-``copies" \(\hat{\nu}_{S_t}\)
    \State \textbf{generate} data from privatized marginals 
    \Statex\hspace{\algorithmicindent}\hspace{\algorithmicindent}\(\{\hat{\nu}_{S_j}\}_{j\leq t}\) to obtain \(D^{(t)}_{DP}\)
\EndFor
\State \Return DP dataset \(D^{(T)}_{DP}\).
\end{algorithmic}
\end{algorithm}
\end{minipage}
\end{figure}

% \begin{minipage}[t]{.45\linewidth}
% \begin{algorithm}[H]
% \caption{Standard data synthesis framework}\label{alg:mainalgo}
% \begin{algorithmic}[1]
% \Require Dataset $D$,  privacy parameters $\epsilon$ and $\delta$
% \State \textbf{select} set $\mathcal S$ of subsets  $S $ of $\{1, \cdots, d\}$
% \State \textbf{privatize}  marginals $\marginal{S}[D]$ to obtain $(\epsilon, \delta)$-DP ``copies" $\measuredmarginal{S}$
% \State \textbf{generate} data from privatized marginals $\measuredmarginal{S}$
% \Statex\Return the DP dataset $ \dpdata$
% \end{algorithmic}
% \end{algorithm}
% \end{minipage}\hfill
% \begin{minipage}{.45\linewidth}
% \begin{algorithm}[H]
% \caption{Sequential query selection}\label{alg:mwem}
% \begin{algorithmic}[1]
% \Require Dataset $D$,  privacy parameters $\epsilon$ and $\delta$, workload $ \mathcal S_W$
% \For{$t = 1, \cdots, T$}
% \State \textbf{select}  $S_t \in  \mathcal S_W$  using $\dpdata^{(t-1)}$ and set 
% %\begin{equation*}
% $\mathcal S_t = \mathcal S_{t-1} \cup \{S_t\}$
% %\end{equation*}
% \State \textbf{run} Step 2 \& 3 in Alg.~\ref{alg:mainalgo} for $\mathcal Q_t$ to obtain $\dpdata^{(t)}$
% \EndFor
% \Statex\Return return DP dataset $ \dpdata^{(T)}$.
% \end{algorithmic}
% \end{algorithm}
% \end{minipage}


% \begin{algorithm}[t!]
% \caption{Standard data synthesis framework}\label{alg:mainalgo}
% \begin{algorithmic}[1]
% \Require Dataset $D$,  privacy parameters $\epsilon$ and $\delta$
% \State \textbf{select} set $\mathcal S$ of subsets  $S $ of $\{1, \cdots, d\}$
% \State \textbf{privatize}  marginals $\marginal{S}[D]$ to obtain $(\epsilon, \delta)$-DP ``copies" $\measuredmarginal{S}$
% \State \textbf{generate} data from privatized marginals $\measuredmarginal{S}$
% \Statex\Return the DP dataset $ \dpdata$
% \end{algorithmic}
% \end{algorithm}

% \begin{algorithm}[t!]
% \caption{Sequential query selection \fy{synthetic ... with adaptive query selection?}}\label{alg:mwem}
% \begin{algorithmic}[1]
% \Require Dataset $D$,  privacy parameters $\epsilon$ and $\delta$, workload $ \mathcal S_W$
% \For{$t = 1, \cdots, T$}
% \State \textbf{select}  $S_t \in  \mathcal S_W$  using $\dpdata^{(t-1)}$ and set 
% %\begin{equation*}
% $\mathcal S_t = \mathcal S_{t-1} \cup \{S_t\}$
% %\end{equation*}

 
% \State \textbf{run} Step 2 \& 3 in Alg.~\ref{alg:mainalgo} for $\mathcal Q_t$ to obtain $\dpdata^{(t)}$
% \EndFor
% \Statex\Return return DP dataset $ \dpdata^{(T)}$.
% \end{algorithmic}
% \end{algorithm}


\paragraph{Marginal privatization}
After a set of marginals has been selected, both frameworks contain a privatization and generation step.
%In the second step, we construct DP-``copies" of the marginals. 
A common choice for privatization is to apply the %standard (doppelt gemoppelt)
Gaussian mechanism \citep{mcsherry2007mechanism}, where we simply add i.i.d.~Gaussian noise to the empirical marginal $\marginal{S}[D]$. The variance $\sigma^2$ of the Gaussian depends on the privacy parameters $\epsilon$ and $\delta$. As a result, we obtain the signed measures  $\measuredmarginal{S}$:
\begin{equation*}
\label{eq:gaussian}
\forall x \in \domain_S:~~
\measuredmarginal{S}(\{x\}) = \marginal{S}[D](\{x\}) + \mathcal N(0, \sigma^2).
\end{equation*}



\paragraph{Data generation} 
Finally, in the last step, we generate a dataset from the noisy estimates of the marginals. 
For this purpose, existing methods 
typically aim to learn a distribution $\hat{p}$ that minimizes the squared loss:
\begin{equation*}
   \sum_{S \in \mathcal S, x \in \domain_S}  \left(\hat p_S(\{x\})  -  \measuredmarginal{S}(\{x\})\right)^2 
\end{equation*}
and then release the private synthetic  data $\dpdata$ by sampling from $\hat p$.
The predominant generation algorithm used by state-of-the-art methods \citep{mst, aim, cai2021data} is PGM~\citep{mckenna2019graphical}, which learns a graphical model using mirror descent. We refer to Section~\ref{sec:relatedwork} for further discussion. 


\subsection{Sliced Wasserstein distance}
\label{sec:slicedwd}
% \fy{In our generation method we crucially utilize the SW ... which we now define}
The optimal transport literature offers a set of natural divergence measures that can be used for particle gradient descent, such as the Sinkhorn divergence or the Wasserstein distance. 
However, the computational complexity of these divergences scales at least quadratically (resp. cubically) with respect to the size of the support of the measures. This clearly defeats the purpose of a computationally efficient generation algorithm that can represent rich datasets with a large number of data points.
Instead, a widely used \citep{wu2019sliced, kolouri2018sliced, deshpande2018generative, deshpande2019max}, computationally efficient alternative is the squared sliced Wasserstein ($\swd_2^2$) distance \citep{bonneel2015sliced}; that is, the averaged squared Euclidean transportation over all $1$-dimensional projections. 

Formally, for $\theta\in\mathbb{R}$, let $g^\theta(x) = \langle x, \theta\rangle$ and $g^\theta_{\#}$ be the pullback measure induced by $g^\theta$. Moreover, let $\lambda^p$ be the uniform distribution over the sphere $\mathcal S^{p-1}$. For any Euclidean subspace $\Omega \subset \mathbb{R}^p$ and probability measures $\mu, \nu \in \mathcal{P}(\Omega)$ with support $\Omega$,  the $q$-SW distance is defined as follows:
\begin{equation}
\label{eq:def_sliced_wasserstein}
    \swd_q(\mu, \nu) =  \left(\int_{\mathcal S^{p -1}} W_q^q(g^\theta_{\#} \mu, g^\theta_{\#}\nu) d\lambda(\theta) \right)^{1/q}, 
\end{equation}
where 
% $p$ is the dimension of $\Omega$ and 
$W_q^q$ is the $q$-th power of the $q$-Wasserstein distance for $q \geq 1$. The Wasserstein distance 
over $1$-dimensional distributions, as it appears in \Cref{eq:def_sliced_wasserstein},
%we can write it using the 
has a well-known closed-form expression:
\begin{equation}
\label{eq:def_wasserstein}
    W_q(g^\theta_{\#} \mu, g^\theta_{\#}\nu)  = \left(\int_u \vert F^{-1}_{g^\theta_{\#} \mu}(u) - F^{-1}_{ g^\theta_{\#}\nu }(u)\vert^q du\right)^{1/q}, 
\end{equation}
 where $F_{g^\theta_{\#} \mu}(u)$ (resp.~$ F_{ g^\theta_{\#}\nu }(u)$) represents the cumulative  function and $F^{-1}$ its inverse. Finally, in the special case of $q=1$, we have the  following identity for the $1$-Wasserstein distance by Vallender \citep{vallander}, which we will later leverage in our algorithm:\begin{equation}\label{eq:vallander}
    W_1(g^\theta_{\#} \mu, g^\theta_{\#}\nu) = \| F_{g^\theta_{\#} \mu}(u) - F_{ g^\theta_{\#}\nu }(u)\|_{L_1(\mathbb R)}.
\end{equation}
Importantly, this identity allows us to naturally extend the definition of the (sliced) $1$-Wasserstein distance to signed measures, as demonstrated by \citet{boedihardjo2022private}.




